\documentclass[11pt]{article}

\usepackage[a4paper,margin=29mm]{geometry}
\usepackage[T1]{fontenc}
\usepackage{lmodern}
\usepackage{microtype}
\usepackage{amsmath,amssymb,amsthm,mathtools,mathrsfs}
\usepackage{array,booktabs}
\usepackage{enumitem}
\usepackage[hidelinks]{hyperref}

\newtheorem{theorem}{Theorem}[section]
\newtheorem{proposition}[theorem]{Proposition}
\newtheorem{lemma}[theorem]{Lemma}
\newtheorem{corollary}[theorem]{Corollary}
\theoremstyle{definition}
\newtheorem{definition}[theorem]{Definition}
\theoremstyle{remark}
\newtheorem{remark}[theorem]{Remark}

\newcommand{\C}{\mathbf C}
\newcommand{\Q}{\mathbf Q}
\newcommand{\Z}{\mathbf Z}
\newcommand{\PP}{\mathbf P}
\newcommand{\Eu}{\mathcal E}
\newcommand{\Gr}{\mathsf{Gr}}
\newcommand{\Nov}{\Lambda}

\title{Owed lemmas for the rank-two atomic residue proof}
\author{Tavis Rudd}
\date{August 2026}

\begin{document}
\maketitle

\noindent\emph{Purpose.} The argument transcribed in
\texttt{2026-08-15-c912-atom-residue-proof.tex} uses, at its first and most
load-bearing step, horizontality of the Poincar\'e pairing for the A-model
connection. Its independent verification
(\texttt{2026-08-15-c912-atom-residue-proof-verification.md}) found that this is
\emph{not} an available import: Katzarkov--Kontsevich--Pantev--Yu state
horizontality only after assuming convergence and passing to the
complex-analytic bundle, whereas the $A$-model $F$-bundle of a projective
variety is non-archimedean analytic. The step is therefore ours to prove. This
note proves it, over the Novikov ring and with no convergence hypothesis, and
lists what else that verification recorded as owed.

\section{Setting}

Let $T$ be a smooth projective variety over $\C$ of dimension $n$, and let
$\Nov$ be the completed Novikov ring of its numerical curve lattice. All
statements below are identities in $\Nov$-modules; no convergence of any
Gromov--Witten generating function is used, and no base change to a
complex-analytic bundle is performed.

Write $H=H^{\mathrm{ev}}(T)\otimes\Nov$ for the even part, $(\,\cdot\,,\cdot\,)$
for the Poincar\'e pairing extended $\Nov$-bilinearly, $\star_\tau$ for the
genus-zero quantum product at a bulk parameter $\tau$, and
\[
  \Gr(\phi)=\tfrac12\bigl(\deg\phi-n\bigr)\phi
  \qquad\text{for homogeneous }\phi,
\]
for the grading operator. The connection is
\[
  \nabla_{u\partial_u}=u\partial_u-u^{-1}\Eu\star_\tau+\Gr,
  \qquad
  \nabla_\xi=\xi+u^{-1}C_\xi,
  \qquad C_\xi=\xi\star_\tau,
  \tag{1}
  \label{eq:connection}
\]
so that horizontal sections satisfy
\[
  u\partial_uy=A(u)y,\quad A(u)=u^{-1}U-\Gr,\ U:=\Eu\star_\tau;
  \qquad
  \xi y=B_\xi(u)y,\quad B_\xi(u)=-u^{-1}C_\xi.
  \tag{2}
  \label{eq:matrices}
\]
The pairing is taken in the usual twisted form: for sections $y,\tilde y$ one
considers $\bigl(y(u),\tilde y(-u)\bigr)$.

\section{The two adjointness identities}

\begin{lemma}[Frobenius, over $\Nov$]
\label{lem:frobenius}
For all $a,b,c\in H$ and every bulk parameter $\tau$,
\[
  (a\star_\tau b,\,c)=(a,\,b\star_\tau c).
\]
In particular every quantum multiplication operator, including $U=\Eu\star_\tau$
and each $C_\xi=\xi\star_\tau$, is self-adjoint for the Poincar\'e pairing.
\end{lemma}

\begin{remark}[The convention this depends on]
That $U$ is a quantum multiplication operator is a convention-sensitive point and
should be cited rather than assumed: in several presentations the Euler operator
on a quantum connection carries an additional derivation term, and if it did
here the lemma would not apply to it. It does not. In
Katzarkov--Kontsevich--Pantev--Yu, (3.31) makes $\Eu_\gamma$ a cohomology class
and (3.30) makes the operator $\mathcal O$-linear multiplication by it, with no
derivation term.
\end{remark}

\begin{proof}
By definition the structure constants of $\star_\tau$ are
\[
  (a\star_\tau b,\,c)
  =\sum_{d}\sum_{m\ge0}\frac{Q^d}{m!}
   \bigl\langle a,b,c,\tau,\dots,\tau\bigr\rangle_{0,m+3,d},
\]
a formal sum over effective numerical classes $d$ with coefficients the
genus-zero Gromov--Witten invariants. Those invariants are defined by
integration against the virtual class of a moduli space of stable maps with
\emph{unordered} marked-point labels up to the specified insertions, hence are
symmetric under permuting their arguments; on the even part no Koszul signs
arise. The displayed expression is therefore totally symmetric in $a,b,c$, and
comparing it with the same expression for $(a,\,b\star_\tau c)$ gives the claim.

Two remarks on what this does and does not use. The identity is an identity of
formal series: it holds coefficientwise in $Q^d$ and in the bulk parameter, so
it is a statement in $\Nov$ and requires no convergence. And it is exactly the
total symmetry of the third derivatives of the genus-zero potential, so it is
the Frobenius property in its usual form, not a strengthening of it.
\end{proof}

\begin{lemma}[The grading operator is anti-self-adjoint]
\label{lem:grading}
For all $a,b\in H$, $\bigl(\Gr a,\,b\bigr)+\bigl(a,\,\Gr b\bigr)=0$.
\end{lemma}

\begin{proof}
Both sides are $\Nov$-bilinear, so it suffices to check on homogeneous classes.
The Poincar\'e pairing of $\phi\in H^k$ with $\psi\in H^\ell$ vanishes unless
$k+\ell=2n$. In that case
\[
  \bigl(\Gr\phi,\psi\bigr)+\bigl(\phi,\Gr\psi\bigr)
  =\Bigl[\tfrac12(k-n)+\tfrac12(\ell-n)\Bigr](\phi,\psi)
  =\tfrac12\bigl[(k+\ell)-2n\bigr](\phi,\psi)=0 .
\]
This is a statement about the cohomological grading and Poincar\'e duality
alone; the Novikov variables play no role, and $\Gr$ is $\Nov$-linear because it
acts on the coefficient ring trivially.
\end{proof}

\begin{remark}
Lemma~\ref{lem:grading} is where the normalization of $\Gr$ by $-n/2$ earns its
keep: with any other additive constant the operator would fail to be
anti-self-adjoint, and horizontality below would acquire a scalar defect.
\end{remark}

\section{Horizontality}

\begin{theorem}[The Poincar\'e pairing is horizontal, over $\Nov$]
\label{thm:horizontal}
Let $y,\tilde y$ be horizontal sections of \eqref{eq:connection}. Then
\[
  \bigl(y(u),\,\tilde y(-u)\bigr)
\]
is constant in $u$ and in every base direction. Equivalently, with $P_0$ the
Gram matrix of the Poincar\'e pairing in a fixed basis,
\[
  A(u)^TP_0+P_0A(-u)=0,
  \qquad
  B_\xi(u)^TP_0+P_0B_\xi(-u)=0 .
  \tag{3}
  \label{eq:horizontality}
\]
\end{theorem}

\begin{proof}
Put $f(u)=y(u)^TP_0\,\tilde y(-u)$. Substituting $w=-u$ shows $u\partial_u$ is
invariant under $u\mapsto-u$, so $(u\partial_u)\bigl(\tilde y(-u)\bigr)
=A(-u)\tilde y(-u)$, and therefore
\[
  u\partial_uf
  =y^T\bigl[A(u)^TP_0+P_0A(-u)\bigr]\tilde y(-u).
\]
Now compute the bracket from \eqref{eq:matrices}:
\[
  \bigl(u^{-1}U-\Gr\bigr)^TP_0+P_0\bigl(-u^{-1}U-\Gr\bigr)
  =u^{-1}\bigl(U^TP_0-P_0U\bigr)-\bigl(\Gr^TP_0+P_0\Gr\bigr).
\]
The first bracket vanishes by Lemma~\ref{lem:frobenius}, since $U$ is a quantum
multiplication operator, and the second by Lemma~\ref{lem:grading}. Hence
$u\partial_uf=0$.

For a base direction, $B_\xi(u)=-u^{-1}C_\xi$ has no $\Gr$ term and is odd in
$u$, so
\[
  B_\xi(u)^TP_0+P_0B_\xi(-u)
  =-u^{-1}\bigl(C_\xi^TP_0-P_0C_\xi\bigr)=0
\]
by Lemma~\ref{lem:frobenius} again, applied to $C_\xi=\xi\star_\tau$. So $f$ is
constant in every base direction as well.
\end{proof}

\begin{corollary}[What transfers to the atomic residue proof]
\label{cor:for-atom}
Equation \textup{(7)} of \texttt{2026-08-15-c912-atom-residue-proof.tex} holds on
the even bundle in the global cohomology frame, with the constant Gram matrix
$P_0$. Its consequences \textup{(8)}, \textup{(10)} and \textup{(11)} follow in
that frame.
\end{corollary}

\begin{proof}
Immediate from Theorem~\ref{thm:horizontal}: with $P(u)=P_0$ constant the
$u\partial_uP$ term of (7) drops and the remainder is
\eqref{eq:horizontality}. The $u^{-1}$ coefficient is $N^TP_0=P_0N$, which is
(8), and the sandwich argument of the target then gives $A_0L\subset L$.
\end{proof}

\begin{remark}[A correction: the pairing is \emph{not} constant on a spectral
factor, and the corollary must not claim it is]
\label{rem:p1-correction}
An earlier version of this corollary asserted additionally that (10) holds
``with $P_1=0$'', on the ground that the pairing is constant. That is
\textbf{false as a general statement} and the error is worth recording, since it
is the kind that hides in bookkeeping.

Theorem~\ref{thm:horizontal} is proved in the global cohomology frame, where the
Gram matrix is the constant $P_0$. The target's equation (7) is applied in the
even \emph{spectral-factor} frame of a single atom, which is reached from the
global frame by a $u$-dependent gauge $F$. In that frame the pairing is
\[
  P(u)=F(u)^TP_0F(-u)
  =P_0|_{\mathrm{block}}-u^2\,\textstyle\sum F_1^TP_0F_1+O(u^3),
\]
which is not constant; $P_1$ vanishes only in a particular
block-off-diagonal Sylvester gauge, so ``$P_1=0$'' is gauge-dependent
bookkeeping and not a fact about the pairing.

Nothing downstream is affected. The target's own sandwich argument derives (8)
and $A_0L\subset L$ with $P_1$ left arbitrary, which is exactly why it was
written that way; the simplification claimed here was never needed. What the
corollary may be cited for is the statement above, in the global frame, and no
more.
\end{remark}

\begin{remark}[What was actually at issue]
The verification's concern was not that horizontality is false but that the
source establishes it only after assuming convergence and passing to the
complex-analytic bundle, while the object the proof works with is
non-archimedean. Theorem~\ref{thm:horizontal} removes that dependence: both
inputs are algebraic identities in $\Nov$-modules, one the total symmetry of
genus-zero Gromov--Witten invariants and the other a statement about the
cohomological grading. Neither refers to a topology on $\Nov$. In particular
nothing here touches the Euler-pairing comparison or the Serre automorphism,
which are the objects the source does defer.
\end{remark}

\section{Still owed}

The same verification recorded four further unstated inputs. They are
bookkeeping rather than substance, but none is written:

\begin{enumerate}[label=(\roman*)]
\item that the even part is a sub-$F$-bundle, so that restricting the whole
  construction to it is legitimate --- note that
  Corollary~\ref{cor:for-atom} does \emph{not} discharge this, and an earlier
  version of that corollary silently consumed it, which was circular;
\item[(i$'$)] the transfer of Theorem~\ref{thm:horizontal} from the global frame
  to a single even spectral factor, carrying a $u$-dependent $P(u)$ as in
  Remark~\ref{rem:p1-correction}, together with spectral orthogonality and
  non-degeneracy of the restricted pairing on that factor. This is the genuine
  residue of the present note and is what the target's (7) actually assumes;
\item uniqueness and gluing of the spectral splitting;
\item the identity principle on irreducible $k$-analytic spaces, which is what
  the crossing-point argument uses when it propagates constancy from a nonempty
  open subset;
\item descent of the invariant through the blowup, projective-bundle and
  disjoint-union equivalences, and not only through the two equivalences in the
  source's own definition.
\end{enumerate}

Separately, and not bookkeeping: the parity ranks $(2,10)$ rest on the source's
assertion that the zero-eigenvalue packet does not split further at the relevant
point, whose cited support carries no proof and covers only a restricted case.
That is the one place where the argument is genuinely conditional, and it is a
finite splitting statement about a single packet rather than an invariance
statement about arbitrary centres.

\end{document}
